% ============================================================================
% GÉNÉRÉ par scripts/exemples-guide.ts — NE PAS ÉDITER À LA MAIN.
% Régénérer : npm run exemples (chaque chiffre est vérifié contre le moteur).
% ============================================================================
\pexemples L'équation~\eqref{eq:rd} appliquée à trois ménages types (montants 2025
en dollars ; chaque composante est la somme des postes détaillés dans les
sections précédentes) :

\begin{center}
\small
\begin{tabular}{l r r r}
\toprule
 & M4 & M6 & M10 \\
\midrule
Revenu de travail ou de retraite & \num{50000} & \num{35000} & \num{20000} \\
Cotisations (postes 1--5) & $-$\num{4622} & $-$\num{2725.81} & $-$\num{762.70} \\
Transferts québécois (postes 6--13) & \num{908.60} & \num{7673.50} & \num{3144.92} \\
Impôt du Québec (poste 19) & $-$\num{3996.40} & $-$\num{1759.24} & $-$\num{100.26} \\
Transferts fédéraux (postes 14--18) & \num{332.30} & \num{8879} & \num{10424.99} \\
Impôt fédéral (poste 19) & $-$\num{3453.30} & $-$\num{0} & $-$\num{197.85} \\
Frais de garde payés (poste 6) & $-$\num{0} & $-$\num{2000} & $-$\num{0} \\
\midrule
\textbf{Revenu disponible} & \textbf{39169.20} & \textbf{45067.45} & \textbf{32509.10} \\
\bottomrule
\end{tabular}
\end{center}
\medskip\noindent À lire en colonne : par exemple, M6 (famille monoparentale,
\D{35000} de travail) \emph{ajoute} \D{7673.50}
de transferts québécois et \D{8879} de
transferts fédéraux à son revenu, pour un revenu disponible de
\D{45067.45} — supérieur à son revenu brut. À l'autre bout,
M4 (personne seule, \D{50000}) cède \D{12071.70}
en prélèvements et ne reçoit que \D{1240.90}
de transferts : revenu disponible de \D{39169.20}.
